\documentclass{article}
\usepackage{readmasters}
\begin{document}

\origpage{293}
\section*{Zur Theorie des eindeutigen Entsprechens algebraischer Gebilde von beliebig vielen Dimensionen.}

\subsection*{Von \textsc{Max Noether} in Göttingen.}

In der Theorie der algebraischen Functionen \emph{einer} complexen Variabeln, die im Wesentlichen auf \textsc{Abel}, \textsc{Puiseux} und \textsc{Riemann} zurückzuführen ist, hat der Letztere eine Classification der Gleichungen aufgestellt${}^{*)}$, welche diese Functionen definiren. Das Characteristische einer Classe ist, dass sich alle derselben angehörigen Gleichungen eindeutig durch rationale Substitutionen in einander transformiren lassen. Wenn $f = 0$ die homogene Gleichung einer Curve $n^{\text{ter}}$ Ordnung mit $d$ Doppel- und $r$ Rückkehrpunkten, so ist${}^{**)}$
\[
p = \frac{n-1 \cdot n-2}{1 \cdot 2} - d - r
\]
die Classenzahl der Classe, zu welcher $f = 0$ gehört, und die Gleichheit dieser Zahl für zwei Curven ist die nothwendige Bedingung für das eindeutige Entsprechen der beiden Curven. Die Zahl $p$, durch welche man die Ordnung des Zusammenhangs derjenigen \textsc{Riemann}'schen Fläche messen kann, in welcher sich die zur Classe $p$ gehörigen algebraischen Functionen geometrisch darstellen lassen, giebt zugleich die Anzahl der allenthalben endlichen Integrale, welche sich auf eine solche Function beziehen.

Eine Erweiterung dieses Theorems auf Functionen zweier Variabeln ist erst vor kurzer Zeit von Herrn \textsc{Clebsch} gegeben worden${}^{***)}$. Wenn sich zwei Flächen rational in einander transformiren lassen, so haben dieselben eine Classenzahl gemeinschaftlich, die Herr \textsc{Clebsch} als das \emph{Geschlecht} der Fläche bezeichnet. Nach dem Theorem, welches Herr \textsc{Clebsch} ohne Beweis mittheilt, ist das Geschlecht $p$ einer Fläche $n^{\text{ter}}$ Ordnung $f = 0$, die nur die gewöhnlichen Singularitäten, Doppel- und Rückkehrcurven, besitzen soll, gleich der Anzahl der in einer Fläche $(n-4)^{\text{ter}}$ Ordnung noch unbestimmt bleibenden Coefficienten, wenn man diese Fläche durch die Doppel- und Rückkehrcurven von $f = 0$ hindurchgehen lässt.

\textbf{*)} Crelle's Journal, Bd.~54.

\textbf{**)} \textsc{Clebsch} und \textsc{Gordan}, Theorie der \textsc{Abel}'schen Functionen.

\textbf{***)} Comptes Rendus de l'Acad. des Sciences vom 21.~Dec.~1868, p.~1238.

\origpage{294}
In einer kürzlich der Societät zu Göttingen mitgetheilten Notiz${}^{*)}$ habe ich diese Theorie nach mehreren Seiten hin ausgedehnt. Der Zweck der folgenden Abhandlung ist nun, das erste Theorem dieser Notiz zu beweisen. Dasselbe ist eine directe Verallgemeinerung des obigen Theorems, und umfasst einmal auch Flächen mit höhern conischen \emph{Knotenpunkten} und höheren vielfachen \emph{Curven}, sodann aber auch die algebraischen Gebilde von \emph{beliebig vielen Dimensionen}, für welche es in dem analogen Umfange, wie dies eben für Flächen angedeutet wurde, eine Classeneintheilung giebt. Zum Beweise werde ich mich einer Methode bedienen, die von den Herren \textsc{Clebsch} und \textsc{Gordan} in ihrem Werke über die „\textsc{Abel}'schen Functionen“ auf die Untersuchung der Functionen \emph{einer} Variabeln angewendet wird, und die den Vortheil besitzt, sich, unter Zufügung einiger Betrachtungen über eindeutiges Entsprechen höherer Gebilde und über algebraische Differentialausdrücke, leicht auf Functionen mehrerer Variabeln ausdehnen zu lassen. Es wird hier nur nöthig, die Beweise für Functionen zweier und dreier Variabeln vollständig durchzuführen, aus denen sich dann das allgemeine Resultat sogleich ergiebt. Ich muss zuerst einige Sätze über das eindeutige Entsprechen zweier Flächen und zweier von drei unabhängig Veränderlichen abhängender Gebilde geben, und gehe dann zur Untersuchung der Differentialausdrücke über, welche sich auf eine algebraische Gleichung beziehen, um hieraus das erwähnte Theorem abzuleiten.

\section*{§ 1.}

\subsection*{Eindeutiges Entsprechen von Flächen.}

Zwei algebraische Gebilde entsprechen sich eindeutig, wenn sie durch rationale Transformationen der Coordinaten in einander übergeführt werden können. Bei diesen eindeutigen Transformationen kommen Eigenthümlichkeiten vor, die eine genauere Betrachtung verlangen.

Die Flächengleichung $n^{\text{ter}}$ Ordnung:
\[
f(x_1,\, x_2,\, x_3,\, x_4) = 0 \tag{1}
\]
werde durch die rationalen Substitutionen
\[
\left\{
\begin{aligned}
\varrho x_1 &= \varphi_1 (y_1,\, y_2,\, y_3,\, y_4) \\
\varrho x_2 &= \varphi_2 (y_1,\, y_2,\, y_3,\, y_4) \\
\varrho x_3 &= \varphi_3 (y_1,\, y_2,\, y_3,\, y_4), \\
\varrho x_4 &= \varphi_4 (y_1,\, y_2,\, y_3,\, y_4)
\end{aligned}
\right. \tag{2}
\]

\textbf{*)} Nachrichten von der Kgl.~Gesellschaft der Wissenschaften zu Göttingen, Jahrgang 1869, No.~15: „Zur Theorie der algebr.~Functionen mehrerer complexen Variabeln.“

\origpage{295}
wo $\varrho$ ein willkürlicher Parameter, die $\varphi$ rationale homogene Functionen $s^{\text{ter}}$ Ordnung ihrer Argumente $y$ sind, in die Form
\[
M \cdot F(y_1,\, y_2,\, y_3,\, y_4) = 0
\]
übergeführt. Der irreductible Factor dieser Gleichung,
\[
F(y_1,\, y_2,\, y_3,\, y_4) = 0, \tag{3}
\]
stellt die transformirte Fläche dar, welche der gegebenen eindeutig entspricht. Ein Factor $M$ muss sich im Allgemeinen absondern, da $F$ nicht eine Function der $\varphi$ sein kann, wenn es möglich sein soll, mit Hülfe der Gleichung $F = 0$ aus den Gleichungen $\varrho x_i = \varphi_i$ die $y$ auch als rationale Functionen der $x$ darzustellen. Sobald $F$ dieser letztern Bedingung genügt, repräsentiren die Gleichungen (2) und (3) eine eindeutige Transformation. Im Allgemeinen entspricht hier jedem Punkte der Fläche $F = 0$ \emph{ein} Punkt $x$ der Fläche $f = 0$, und umgekehrt; im Besondern kann es jedoch eintreten, dass für einzelne oder auch für eine einfach unendliche Reihe von Werthsystemen der $y$ die Functionen $\varphi$ sämmtlich verschwinden, also die $x$ aus den obigen Substitutionsformeln unbestimmt hervorgehen, sowie in den umgekehrten Formeln die $y$ unbestimmt werden können. Die Functionen $\varphi$ lassen sich als Gleichungen von Flächen betrachten, die mit $F = 0$ Punkte oder eine Curve gemein haben können. Für solche Werthsysteme der $y$ ergeben sich die entsprechenden Werthsysteme der $x$, indem man das Werthsystem der $y$ sich dem betrachteten Punkte $y$ annähern lässt. Man kann dabei eines der $y$, z.~B. $y_4$, als Constante ansehen, dann für $y_1$, $y_2$, $y_3$ resp. $y_1 + \varepsilon \eta_1$, $y_2 + \varepsilon \eta_2$, $y_3 + \varepsilon \eta_3$ setzen, die $\varphi$ nach aufsteigenden Potenzen von $\varepsilon$ ordnen und diese Grösse gegen $0$ convergiren lassen. Wenn die Functionen $\varphi$ in diesem Punkte $y$ einfach verschwinden, so hebt sich dabei der Factor $\varepsilon$ heraus, den man in $\varrho$ eingehen lässt; indem man hierauf $\varepsilon = 0$ setzt, gehen die Gleichungen (2) in die folgenden über:
\[
\left\{
\begin{aligned}
\varrho x_1 &= \varphi'_1 (y_1) \cdot \eta_1 + \varphi'_1 (y_2) \cdot \eta_2 + \varphi'_1 (y_3) \cdot \eta_3 \\
\varrho x_2 &= \varphi'_2 (y_1) \cdot \eta_1 + \varphi'_2 (y_2) \cdot \eta_2 + \varphi'_2 (y_3) \cdot \eta_3 \\
\varrho x_3 &= \varphi'_3 (y_1) \cdot \eta_1 + \varphi'_3 (y_2) \cdot \eta_2 + \varphi'_3 (y_3) \cdot \eta_3 \\
\varrho x_4 &= \varphi'_4 (y_1) \cdot \eta_1 + \varphi'_4 (y_2) \cdot \eta_2 + \varphi'_4 (y_3) \cdot \eta_3
\end{aligned}
\right. . \tag{4}
\]
Zugleich hat man, wenn $F = 0$ in $y$ ebenfalls einen einfachen Punkt besitzt, für die $\eta$ die Beziehung:
\[
0 = F'(y_1) \cdot \eta_1 + F'(y_2) \cdot \eta_2 + F'(y_3) \cdot \eta_3 . \tag{5}
\]
Die Elimination von $\varrho$ und der $\eta$ aus diesen Gleichungen (4) und (5) führt auf die Gleichungen zweier Ebenen, deren Schnittlinie auf der Fläche $f = 0$ liegt und dem betrachteten Punkte $y$ von $F = 0$ entspricht. Dies giebt den Satz:

\origpage{296}
„Wenn für einen einfachen Punkt von $F = 0$ die Functionen $\varphi$ in (2) einfach verschwinden, so entspricht diesem Punkte eine Gerade auf $f = 0$.“

Im Falle die $\varphi$ in dem einfachen Punkte von $F = 0$ sämmtlich $\mu$fach verschwinden, hebt sich in der Entwicklung der Gleichungen (2) nach Potenzen von $\varepsilon$ der Factor $\varepsilon^\mu$ heraus, indem alle Glieder bis zur $\mu^{\text{ten}}$ Ordnung verschwinden, und diese Gleichungen nehmen die Gestalt an:
\[
\left\{
\begin{aligned}
\varrho x_i &= \frac{\partial^\mu \varphi_i}{\partial y_1^\mu} \cdot \eta_1^\mu + \mu \cdot \frac{\partial^\mu \varphi_i}{\partial y_1^{\mu-1} \partial y_2} \cdot \eta_1^{\mu-1} \cdot \eta_2 + \mu \cdot \frac{\partial^\mu \varphi_i}{\partial y_1^{\mu-1} \partial y_3} \cdot \eta_1^{\mu-1} \cdot \eta_3 \\
&\qquad\quad + \dots + \frac{\partial^\mu \varphi_i}{\partial y_3^\mu} \cdot \eta_3^\mu
\end{aligned}
\right. \tag{6}
\]
\[
(i = 1,\, 2,\, 3,\, 4).
\]
Indem man eines der $\eta$ aus Gleichung (5) linear und homogen durch die beiden andern ausdrückt und in die Gleichungen (6) einsetzt, findet man die Coordinaten $x$ der dem Punkte $y$ entsprechenden Curve als rationale homogene Functionen $\mu^{\text{ter}}$ Ordnung zweier Parameter $\eta$, was zu dem allgemeineren Satze führt:

„Wenn für einen einfachen Punkt von $F = 0$ die Functionen $\varphi$ in (2) $\mu$fach verschwinden, so entspricht diesem Punkte auf der Fläche $f = 0$ eine Curve $\mu^{\text{ter}}$ Ordnung, welche die Eigenschaft hat, ihre Coordinaten als rationale Functionen eines Parameters ausdrücken zu lassen.“

Dass diese letztere Eigenschaft stattfinden muss, liess sich schon unmittelbar aus dem obigen specielleren Satze erkennen. Denn man kann immer an Stelle derjenigen Transformation, welche von der Fläche $f = 0$ mit der Curve $\mu^{\text{ter}}$ Ordnung zur Fläche $F = 0$ führt, zwei successive eindeutige Transformationen setzen, von denen die erste von $f = 0$ zu einer Fläche $f_1 = 0$ mit einer Geraden führt, die zweite von $f_1 = 0$ zu $F = 0$. Wenn dann die Curve $\mu^{\text{ter}}$ Ordnung und die Gerade \emph{einem} Punkte $y$ von $F = 0$ entsprechen, so müssen, da die beiden Flächen $f$ und $f_1$ sich eindeutig entsprechen, auch die beiden Curven sich Punkt für Punkt entsprechen und somit das gleiche Geschlecht $p$ haben. Die Transformation zwischen $f$ und $f_1$ findet man, indem man die Ausdrücke der $y$ aus Transformationsformeln zwischen $f_1$ und $F$, die sich immer angeben lassen, in die Formeln zwischen $f$ und $F$ einsetzt. Von diesem \emph{Princip der successiven Transformation} werden wir noch weiterhin zur Vereinfachung der Untersuchung Gebrauch machen.

Einem $\nu$fachen Knotenpunkte von $F = 0$ entspricht auf $f = 0$ eine Curve, welche durch die Gleichung:

\origpage{297}
\[
0 = \frac{\partial^\nu F}{\partial y_1^\nu} \cdot \eta_1^\nu + \nu \frac{\partial^\nu F}{\partial y_1^{\nu-1} \partial y_2} \eta_1^{\nu-1} \eta_2 + \nu \frac{\partial^\nu F}{\partial y_1^{\nu-1} \partial y_3} \eta_1^{\nu-1} \cdot \eta_3 + \dots + \frac{\partial^\nu F}{\partial y_3^\nu} \cdot \eta_3^\nu \tag{7}
\]
in Verbindung mit den Gleichungen (4) oder (6) bestimmt wird. Diese Gleichungen führen zu dem folgenden Satze:

Wenn für einen $\nu$fachen Knotenpunkt von $F = 0$ die Functionen $\varphi$ sämmtlich $\mu$fach verschwinden, so entspricht diesem Knotenpunkte auf $f = 0$ eine Curve, deren Geschlecht durch Gleichung (7) bestimmt ist. Den $\nu$ Fortschreitungsrichtungen innerhalb einer durch den Knotenpunkt gehenden Ebene entsprechen $\nu$ Punkte dieser Curve, die auf einer Curve $\mu^{\text{ter}}$ Ordnung liegen, welche die Eigenschaft hat, ihre Coordinaten rational durch einen Parameter ausdrücken zu lassen.

Noch einen weitern ganz ähnlichen Satze erhält man für die Curve, welche eine vielfache Curve von $F = 0$ abbildet. Wenn $F$ die Form hat:
\[
0 = \varkappa \cdot A^\nu + \varkappa_1 A^{\nu-1} \cdot B + \dots + \varkappa_\nu B^\nu = 0, \tag{8}
\]
und die Functionen $\varphi$ durch die Curve $A = 0$, $B = 0$ einfach hindurchgehen, also die Substitutionsformeln die Form annehmen:
\[
\varrho x_i = \lambda_i A + \mu_i B, \tag{9}
\]
wo die $\varkappa$, $\lambda$, $\mu$ Functionen der $y$ sind, so findet man die $\nu$ verschiedenen Punkte $x$, welche einem gegebenen Punkte $y$ der vielfachen Curve $A = 0$, $B = 0$ entsprechen, indem man die Coordinaten dieses Punktes in die $\varkappa$, $\lambda$, $\mu$, und sodann die $\nu$ Werthe von $\frac{A}{B}$ aus (8) in (9) einsetzt. Die Coordinaten dieser Punkte $x$ werden dann von der Form:
\[
\sigma x_i = \lambda_i A_s + \mu_i B_s, \quad (s = 1,\, 2,\, \dots\, \nu),
\]
woraus folgt, dass die $\nu$ Punkte $x$ auf einer Geraden liegen. Allgemeiner hat man:

Wenn durch eine $\nu$fache Curve von $F = 0$ die Functionen $\varphi$ $\mu$fach hindurchgehen, so entsprechen einem Punkte dieser Curve $\nu$ Punkte auf $f = 0$, die auf einer Curve $\mu^{\text{ter}}$ Ordnung vom Geschlechte $0$ liegen.

Die Abbildung von vielfachen Curven ergiebt sich, indem man aus den Gleichungen:
\[
\left\{
\begin{aligned}
0 &= \varkappa \cdot \eta_1^\nu + \varkappa_1 \eta_1^{\nu-1} \cdot \eta_2 + \dots + \varkappa_\nu \cdot \eta_2^\nu, \\
A &= 0, \quad B = 0, \\
\varrho x_i &= \lambda_i \eta_1 + \mu_i \eta_2,
\end{aligned}
\right.
\]
die Verhältnisse der $y$, ferner $\frac{\eta_1}{\eta_2}$ und $\varrho$ eliminirt.

Ich will noch bemerken, dass Rückkehrpunkte von $\nu$fachen Curven oder eine Rückkehrcurve selbst ebenfalls in diesen Satz einge-

\origpage{298}
schlossen sind, und nur ein unendlich nahes Zusammenrücken der betreffenden Punkte $x$ bedingen.

Wenn in den bisher betrachteten Fällen die Functionen $\varphi$ in einem singulären Punkte $y$ von $F = 0$ nicht verschwinden, so bildet sich der Punkt $y$ durch einen einzelnen Punkt $x$ auf $f = 0$ ab, und es ergiebt sich zu jeder Fortschreitungsrichtung in $y$ auf $F = 0$ eine solche in dem entsprechenden Punkte $x$ von $f = 0$, so dass die Curven auf $F$, die in $y$ einen $\nu$fachen Punkt zeigen, einen ebensolchen in ihrer Abbildung in $x$ zeigen müssen. Wenn endlich drei der Functionen $\varphi$ und die Fläche $F = 0$ mehrere Punkte gemein haben, so entspricht diesen Punkten im Allgemeinen \emph{ein} Knotenpunkt auf $f = 0$, der jedoch eine höhere Singularität zeigt, als die eines conischen Punktes, ein Fall, den wir daher aus unseren späteren Betrachtungen ausschliessen werden. So können einem biplanaren Punkte, der sich im Allgemeinen durch zwei sich schneidende gerade Linien oder diesen entsprechende Curven abbildet, speciell zwei getrennte Punkte entsprechen, wie z.~B. in der Reciprokalfläche.

Diese und alle weiteren Fälle ergeben sich, wenn man mehrere Transformationen nach einander ausführt und bei diesen nur die bisher entwickelten Abbildungen eintreten lässt.

\section*{§ 2.}

\subsection*{Eindeutiges Entsprechen von höheren Mannigfaltigkeiten.}

Eine homogene Gleichung $n^{\text{ten}}$ Grades mit $r + 2$ Variabeln:
\[
f(x_1,\, x_2,\, \dots,\, x_{r+2}) = 0
\]
definirt das Verhältniss zweier der Variabeln als algebraische Function von $r$ unabhängig veränderlichen complexen Grössen und stellt eine $2r$fach unendliche Mannigfaltigkeit vor. Zur Abkürzung werde ich im Folgenden eine $2h$fach unendliche Mannigfaltigkeit als eine $\infty^{2.h}$ bezeichnen. Die auf $f = 0$ liegenden Mannigfaltigkeiten sind dann als $\infty^{2.r-1}$, $\infty^{2.r-2}$, \dots\ $\infty^{2.1}$ (Curven), $\infty^0$ (Punkte) zu bezeichnen.

Wenn sich zwei $\infty^{2.r}$, $f(x) = 0$ und $F(y) = 0$, Punkt für Punkt eindeutig entsprechen, so wird es im Besondern eintreten, je nachdem sämmtliche oder einige der Functionen $\varphi$ der eindeutigen Transformation mehrfach verschwinden, dass einer $\infty^{2.h}$ auf $F = 0$ eine $\infty^{2.\varkappa}$ auf $f = 0$ entspricht, wo $\varkappa \geqq h$, $r > h \geqq 0$.

Da bereits der Fall einer $\infty^{2.3}$ die bei den Abbildungen eintretenden Verhältnisse deutlich übersehen lässt, so beschränke ich mich hier auf diesen Fall, der überdies noch dadurch weiteres geometrisches Interesse hat, dass auf ihn die Gleichungen der \emph{Liniencomplexe} zurückführen.

Die Gleichungen (2) und (3) des vorhergehenden Paragraphen stellen wieder, wenn sie zu den folgenden\ednote{In the following equation the fourth argument of $F$ is printed $y_1$ instead of $y_4$; an apparent misprint, as the preceding line shows.}

\origpage{299}
\[
\begin{gathered}
\varrho x_i = \varphi_i (y_1,\, y_2,\, y_3,\, y_4,\, y_5), \quad (i = 1,\, 2,\, \dots\, 5) \\
0 = F(y_1,\, y_2,\, y_3,\, y_1,\, y_5)
\end{gathered}
\]
erweitert werden, eine eindeutige Transformation dar, durch welche man von der $\infty^{2.3}$, $F = 0$, zu der $\infty^{2.3}$
\[
0 = f(x_1,\, x_2,\, x_3,\, x_4,\, x_5),
\]
und umgekehrt, übergeht.

Für diejenigen Werthsysteme der $y$, für welche die $\varphi_i$, und zwar nicht sämmtlich mehrfach, verschwinden, ergeben sich die zugehörigen Werthsysteme der $x$ aus den Gleichungen:\ednote{In the following display the last term is printed with $\eta_1$ instead of $\eta_4$, matching the argument $y_4$. An apparent misprint, reproduced here as printed.}
\[
\varrho x_i = \varphi'_i(y_1) \cdot \eta_1 + \varphi'_i(y_2) \cdot \eta_2 + \varphi'_i(y_3) \cdot \eta_3 + \varphi'_i(y_4) \cdot \eta_1
\]
\[
(i = 1,\, 2,\, \dots\, 5).
\]
Im Falle hier für einen einfachen Punkt $y$ von $F = 0$ die Functionen $\varphi$ sämmtlich einfach, oder eines der $\varphi$ oder mit Hülfe von $F = 0$ zwei der $\varphi$ zweifach, die übrigen einfach verschwinden, folgt, dass dem Punkte $y$ eine auf $f = 0$ liegende \emph{Ebene} entspricht.

Kann man mit Hülfe von $F = 0$ in einem Punkte $y$ drei lineare Combinationen der $\varphi$ im $2^{\text{ten}}$ Grade verschwinden machen, während die übrigen $\varphi$ für diesen Punkt nur einfach verschwinden, so ergiebt die Elimination der $\varrho$ und $\eta$ drei lineare Gleichungen, deren Lösungen $f = 0$ genügen, also eine auf $f = 0$ liegende \emph{Gerade}, welche dem Punkte $y$ entspricht.

Wendet man nun auf den Fall des mehrfachen Verschwindens der $\varphi$ das Princip der successiven Transformation an, so hat man den folgenden Satz:

Im Falle die Functionen $\varphi$ in einem einfachen Punkte von $F = 0$ sämmtlich verschwinden, entspricht diesem Punkte auf $f = 0$ eine Fläche, die sich rational durch zwei Parameter, oder eine Curve, die sich rational durch einen Parameter ausdrücken lässt.

Von grösserem Interesse ist der Fall, dass einer Curve auf $F = 0$ eine Fläche auf $f = 0$ entspricht. Wenn $A = 0$, $B = 0$, $C = 0$ die Gleichungen der Curve sind, so werden hier die Transformationsformeln für den Fall, dass die $\varphi$ nur einfach verschwinden:
\[
\begin{gathered}
0 = F(y) = \lambda A + \mu B + \nu C \\
\varrho x_i = \lambda_i A + \mu_i B + \nu_i C \quad (i = 1,\, 2,\, \dots\, 5),
\end{gathered}
\]
und die Elimination der $A$, $B$, $C$ führt zu dem unvollständigen Determinantensystem:
\[
0 =
\begin{vmatrix}
x_1 & \lambda_1 & \mu_1 & \nu_1 \\
x_2 & \lambda_2 & \mu_2 & \nu_2 \\
x_3 & \lambda_3 & \mu_3 & \nu_3 \\
x_4 & \lambda_4 & \mu_4 & \nu_4 \\
x_5 & \lambda_5 & \mu_5 & \nu_5 \\
0 & \lambda & \mu & \nu
\end{vmatrix} .
\]
Denkt man sich nun in die $\lambda$, $\mu$, $\nu$ die Coordinaten eines Punktes $y$

\origpage{300}
der Curve $A = 0$, $B = 0$, $C = 0$ eingesetzt, so stellt dieses System die dem Punkte $y$ entsprechende Curve, eine \emph{Gerade} vor.

Die Gleichungen $A = 0$, $B = 0$, $C = 0$ bestimmen aber eine Curve, der eine ganze Fläche entspricht. Diese erhält man, wenn man aus den Gleichungen:
\[
\left\{
\begin{aligned}
0 &= \lambda \eta_1 + \mu \eta_2 + \nu \eta_3 \\
\varrho x_i &= \lambda_i \eta_1 + \mu_i \eta_2 + \nu_i \eta_3, \quad (i = 1,\, 2,\, \dots\, 5) \\
A &= 0, \quad B = 0, \quad C = 0
\end{aligned}
\right.
\]
drei der Verhältnisse der $y$ und eines der Verhältnisse $\frac{\eta_1}{\eta_2}$, $\frac{\eta_1}{\eta_3}$ eliminirt. Die $x$ erscheinen dann als rationale lineare homogene Functionen zweier der $\eta$, deren Coefficienten selbst algebraische Functionen einer Variabeln sind, deren Irrationalität gegeben ist durch die Gleichungen $A = 0$, $B = 0$, $C = 0$. Der Curve auf $F = 0$ entspricht also auf $f = 0$ eine aus geraden Linien gebildete Fläche, eine \emph{Regelfläche}. Die in den Gleichungen $A = 0$, $B = 0$, $C = 0$ enthaltene Irrationalität kann man durch Einführen zweier neuen Variabeln $s$, $z$ auf die in einer Gleichung $\psi(s,\, z) = 0$ enthaltene zurückführen, die, wenn sie von der $m^{\text{ten}}$ Ordnung wird, einen Kegel $m^{\text{ter}}$ Ordnung, also eine Fläche $m^{\text{ter}}$ Ordnung mit einem $m$fachen Knotenpunkte darstellt. Mit Hülfe des unten folgenden Theorems, dessen Beweis den Gegenstand des vorliegenden Aufsatzes bildet, ergiebt sich aus dieser Bemerkung der wichtige Satz, dass die einer einfachen Curve auf $F = 0$ entsprechenden Regelflächen auf $f = 0$ das Flächengeschlecht $0$ haben. Nimmt man noch die leicht zu beweisende Bemerkung hinzu, dass die Kegel sich im Allgemeinen nicht eindeutig auf einer Ebene abbilden lassen, so findet man mittelst Erweiterung durch das Princip der successiven Transformation den Satz:

Im Falle die Functionen $\varphi$ längs einer einfachen Curve von $F = 0$ sämmtlich $\mu$fach verschwinden, entspricht jedem Punkte der Curve eine Curve $\mu^{\text{ter}}$ Ordnung auf $f = 0$, die sich rational durch einen Parameter ausdrücken lässt. Die der Curve auf $F$ entsprechende Fläche auf $f$ wird durch die Bewegung dieser veränderlichen rationalen Curve $\mu^{\text{ter}}$ Ordnung erzeugt und hat das Geschlecht $0$, lässt sich aber im Allgemeinen nicht rational durch zwei Parameter ausdrücken.

Von den noch übrigen Besonderheiten, welche bei Abbildungen zweier $\infty^{2.3}$ auf einander eintreten, will ich nur noch die folgenden erwähnen. Wenn $F = 0$ einen Knotenpunkt enthält und die $\varphi$ verschwinden in diesem Punkte, so löst sich derselbe auf $f = 0$ im Allgemeinen in eine Fläche auf, deren Irrationalität mit derjenigen übereinstimmt, welche die Gleichung $F = 0$ in der Nähe des Knotenpunktes besitzt. Eine vielfache Curve von $F = 0$, längs deren die $\varphi$

\origpage{301}
verschwinden, löst sich in ihrer Abbildung auf $f = 0$ ebenfalls in eine Fläche auf. Bei $\mu$fachem Verschwinden der $\varphi$ längs dieser Curve von $F$ entspricht jedem Punkte dieser vielfachen Curve eine Curve auf $f = 0$ (von höherem Geschlecht im Allgemeinen), die auf einer Fläche vom Grade $\mu$ liegt, deren Coordinaten sich rational durch zwei Parameter ausdrücken lassen. Auch eine $\nu$fache Fläche von $F = 0$ löst sich, wenn die $\varphi$ in allen ihren Punkten verschwinden, in ihrer Abbildung auf $f = 0$ in eine Fläche auf, und jedem ihrer Punkte entsprechen $\nu$ Punkte auf $f$, die auf einer Curve vom Geschlecht $0$ liegen.

Auch hier gilt die Bemerkung, dass sich alle weiteren Fälle mittelst des Princips der successiven Transformation erledigen lassen. Man führt mittelst desselben alle Fälle auf diejenigen des Entsprechens zwischen Fläche und Punkt und zwischen Fläche und Curve zurück.

\section*{§ 3.}

\subsection*{Die algebraischen Differentialausdrücke und deren Integrale.}

Ich betrachte diejenigen Irrationalitäten, welche mit der Gleichung einer irreductibeln $\infty^{2.r}$:
\[
f(s,\, z_1,\, z_2,\, \dots\, z_r) = 0,
\]
die $s$ als algebraische Function von $r$ complexen Variabeln $z$ definirt, verbunden sind. Zur symmetrischen Darstellung der Differentialausdrücke, welche sich auf $f = 0$ beziehen, denke ich mir diese Gleichung in homogenen Variabeln geschrieben in der Form:
\[
f(x_1,\, x_2,\, \dots\, x_{r+2}) = 0, \tag{1}
\]
und führe hier $r$ neue unabhängig Veränderliche derart ein, dass hierdurch keine weitere Irrationalität in die Gleichung eingeht. Man erreicht dies dadurch, dass man als neue unabhängig Veränderliche $r$ rationale Functionen, $\lambda_1$, $\lambda_2$, \dots\ $\lambda_r$, von $s$, $z_1$, \dots\ $z_r$ (identisch mit homogenen rationalen Functionen $0^{\text{ter}}$ Ordnung von $x_1$, $x_2$, \dots\ $x_{r+2}$) einführt; denn diese Substitution geschieht durch Gleichungen von der Form:
\[
\varphi_1 + \lambda_1 \varphi'_1 = 0, \quad \varphi_2 + \lambda_2 \varphi'_2 = 0, \ \dots\ \varphi_r + \lambda_r \varphi'_r = 0, \quad s = s,
\]
die als die Formeln einer eindeutigen Transformation aufzufassen sind, bei der sich, mit Hülfe von $f = 0$, auch die $x$ rational durch
\[
s,\, \lambda_1,\, \lambda_2,\, \dots\, \lambda_r
\]
ausdrücken. Ich denke mir diese Formeln zur Definition der $x$ als \emph{algebraische Functionen der $r$ complexen Variabeln $\lambda$} benutzt.

Es sei Gleichung (1) homogen von der $n^{\text{ten}}$ Ordnung. Man hat für alle Werthe der unabhängig Veränderlichen $\lambda$ die $r+1$ Gleichungen:

\origpage{302}
\begin{align*}
\frac{\partial f}{\partial x_1} \cdot x_1 + \frac{\partial f}{\partial x_2} \cdot x_2 + \dots + \frac{\partial f}{\partial x_{r+2}} \cdot x_{r+2} &= nf = 0 \\
\frac{\partial f}{\partial x_1} \cdot \frac{\partial x_1}{\partial \lambda_1} + \frac{\partial f}{\partial x_2} \cdot \frac{\partial x_2}{\partial \lambda_1} + \dots + \frac{\partial f}{\partial x_{r+2}} \cdot \frac{\partial x_{r+2}}{\partial \lambda_1} &= 0 \\
\frac{\partial f}{\partial x_1} \cdot \frac{\partial x_1}{\partial \lambda_2} + \frac{\partial f}{\partial x_2} \cdot \frac{\partial x_2}{\partial \lambda_2} + \dots + \frac{\partial f}{\partial x_{r+2}} \cdot \frac{\partial x_{r+2}}{\partial \lambda_2} &= 0 \\
\dots\dots\dots\dots\dots\dots\dots\dots\dots\dots\dots\dots & \dots \\
\frac{\partial f}{\partial x_1} \cdot \frac{\partial x_1}{\partial \lambda_r} + \frac{\partial f}{\partial x_2} \cdot \frac{\partial x_2}{\partial \lambda_r} + \dots + \frac{\partial f}{\partial x_{r+2}} \cdot \frac{\partial x_{r+2}}{\partial \lambda_r} &= 0.
\end{align*}
Hieraus folgt, dass sich die $\frac{\partial f}{\partial x_i}$ verhalten wie die Unterdeterminanten nach den $c_i$ der Determinante:
\[
\begin{vmatrix}
c_1 & c_2 & \cdot & \cdot & c_{r+2} \\
x_1 & x_2 & \cdot & \cdot & x_{r+2} \\
\frac{\partial x_1}{\partial \lambda_1} & \frac{\partial x_2}{\partial \lambda_1} & \cdot & \cdot & \frac{\partial x_{r+2}}{\partial \lambda_1} \\
\frac{\partial x_1}{\partial \lambda_2} & \frac{\partial x_2}{\partial \lambda_2} & \cdot & \cdot & \frac{\partial x_{r+2}}{\partial \lambda_2} \\
\cdot & \cdot & \cdot & \cdot & \cdot \\
\frac{\partial x_1}{\partial \lambda_r} & \frac{\partial x_2}{\partial \lambda_r} & \cdot & \cdot & \frac{\partial x_{r+2}}{\partial \lambda_r}
\end{vmatrix},
\]
und somit ist der Ausdruck
\[
\frac{\Sigma \pm c_1\, x_2\, \frac{\partial x_3}{\partial \lambda_1} \frac{\partial x_4}{\partial \lambda_2} \dots \frac{\partial x_{r+2}}{\partial \lambda_r}}{\displaystyle\sum_i c_i\, \frac{\partial f}{\partial x_i}},
\]
der in Bezug auf die $x$ symmetrisch ist, unabhängig von den ganz willkürlichen Grössen $c$.

Bezeichnet ferner $\Theta$ eine rationale homogene Function $n - r - 2^{\text{ter}}$ Ordnung der $x$, so ist
\[
\Omega = \frac{\Theta \cdot \Sigma \pm c_1\, x_2\, \frac{\partial x_3}{\partial \lambda_1} \frac{\partial x_4}{\partial \lambda_2} \dots \frac{\partial x_{r+2}}{\partial \lambda_r}}{\displaystyle\sum_i c_i\, \frac{\partial f}{\partial x_i}}
\]
ein nach den $\lambda$ genommener Differentialausdruck, der, als homogene Function $0^{\text{ter}}$ Ordnung, nur noch von den Verhältnissen der $x$ abhängt, daher ein algebraischer Differentialausdruck in Bezug auf die $\lambda$, dessen Irrationalität jedoch nur durch Gleichung (1) bedingt ist.

In der Theorie der Functionen \emph{einer} Variabeln betrachtet man allenthalben endliche einfache Integrale, genommen über Differentialausdrücke, die $\Omega$ analog sind, und weist nach, dass bei einer ein-

\origpage{303}
deutigen Transformation die Eigenschaft der Endlichkeit erhalten bleibt. Um hier einen ähnlichen Gang einzuschlagen und Anhaltspunkte für eine geeignete Bestimmung der Function $\Theta$ im Zähler von $\Omega$ zu gewinnen, wollen wir einen Blick auf die vielfachen Integrale von $\Omega$ werfen:
\[
\omega = \int \frac{\Theta \cdot \Sigma \pm c_1\, x_2\, \frac{\partial x_3}{\partial \lambda_1} \frac{\partial x_4}{\partial \lambda_2} \dots \frac{\partial x_{r+2}}{\partial \lambda_r}}{\displaystyle\sum_i c_i\, \frac{\partial f}{\partial x_i}} \cdot d\lambda_1\, d\lambda_2 \dots d\lambda_r .
\]
Das Werthgebiet, über welches das $r$fache Integral ausgedehnt werden soll, kann man sich dadurch festgelegt denken, dass man der Reihe nach die $\lambda$ festbestimmte, einfach unendliche Werthreihen durchlaufen lässt, von constanten Grenzen an bis zu einem Werthsysteme $\lambda$, und zugleich jedem Werthsysteme $\lambda$ bestimmte Werthe der $x$ zuertheilt, wodurch dann $\omega$ eine bestimmte Function der $\lambda$ wird, wenn innerhalb des Integrationsgebiets keine Unstetigkeiten vorkommen. Eine Aenderung der Integrationsweise innerhalb desselben Integrationsgebiets, d.~h. innerhalb einer gegebenen $r$fach unendlichen Werthreihe der $x$, kann durch Einführung anderer Functionen $\mu$ statt der $\lambda$ bewirkt werden, und geschieht dann durch eine eindeutige Transformation, hat also auf die Irrationalität, die sich in den Werthen der Integrale $\omega$ ausprägt, keinen Einfluss. Auch in Bezug auf Unstetigkeiten verhält sich das Integral $\omega$ bei veränderter Integrationsweise, wie das Integral eines eindeutig transformirten Ausdrucks $\Omega'$, den wir hier aber so bestimmen wollen, dass er ganz analoge Unstetigkeiten besitzen soll, wie $\Omega$ selbst.

Einige einfache Betrachtungen über die Unstetigkeiten der vielfachen Integrale ergeben nun Folgendes:

Vor Allem wird $\omega$ unstetig in denjenigen Werthsystemen der $\lambda$, für welche der Nenner von $\Theta$ verschwindet. Um nur endliche Integrale zu betrachten, setzen wir daher voraus, dass $\Theta$ eine \emph{ganze}, homogene Function $(n - r - 2)^{\text{ter}}$ Ordnung der $x$ sei.

Dann sind die Unstetigkeiten von $\omega$ auf diejenigen Stellen beschränkt, in denen
\[
\sum_i c_i\, \frac{\partial f}{\partial x_i}
\]
verschwindet. Die Punkte, in welchen
\[
\sum_i c_i\, \frac{\partial f}{\partial x_i} = 0,
\]
ohne dass die $\frac{\partial f}{\partial x_i}$ sämmtlich einzeln verschwinden, repräsentiren zusammen das \emph{Verzweigungsgebilde} von $f = 0$, das hier eine $\infty^{2(r-1)}$ ist.

\origpage{304}
Für Flächen fällt dieses Gebilde mit der Curve zusammen, in welcher $f = 0$ vom Tangentenkegel des Punktes $c$ berührt wird. Irgend ein $r$fach unendliches Integrationsgebiet hat im Allgemeinen mit dem Verzweigungsgebilde wenigstens eine $\infty^{r-2}$ gemein, so dass man genöthigt ist, über diese Gebilde hinweg zu integriren. Indessen ergiebt sich, dass die Unstetigkeit der Differentialausdrücke in den Verzweigungsgebilden bei den Integralen wegfällt.

Den Punkten, in welchen die $\frac{\partial f}{\partial x_i}$ sämmtlich verschwinden, entsprechen die auf $f = 0$ liegenden vielfachen Mannigfaltigkeiten. Wir werden hier nur diejenige Eigenschaft der Function $\Theta$ anführen, welche man durch einfache Betrachtungen für die Endlichkeit der Integrale als \emph{nothwendig} erkennt, und von der wir im Folgenden allein Gebrauch zu machen haben. Wir setzen dabei voraus, dass die vielfachen, auf $f = 0$ liegenden Mannigfaltigkeiten nicht selbst wieder specielle Singularitäten enthalten, die denen analog wären, welche bei einer Fläche durch das Zusammenfallen von Tangentenebenen längs einer vielfachen Curve oder durch den Uebergang eines conischen Knotenpunkts in einen planaren eintreten, und die sich in der, für die Punkte in der Nähe der betreffenden Singularität modificirten Gleichung $f = 0$ durch das Absondern von Factoren oder auch, bei vielfachen $\infty^{2(r-1)}$ auf $f = 0$, durch das Zusammenfallen von Factoren ausdrücken würden. Diese Singularitäten erfordern jedesmal ganz specielle Untersuchungen. Nur der Fall, dass längs einer doppelt zu zählenden $\infty^{2(r-1)}$ auf $f = 0$ die beiden linearen Factoren, in welche $f = 0$ in der Nähe jedes Punktes dieser $\infty^{2(r-1)}$ zerfällt, zusammenfallen, also speciell der Fall einer Rückkehrcurve auf einer Fläche $f = 0$, ist hier noch mit berücksichtigt.

Die Bedingung, dass die Integrale, genommen über allgemeine Integrationsgebiete, die durch Werthsysteme von vielfachen $\infty^{2h}$ hindurchgehen, nicht unendlich werden sollen, erfordert, dass $\Theta$ längs jeder $\mu$fach zählenden $\infty^{2h}$ auf $f = 0$ $(\mu - r + h)$mal verschwinde, d.~h. diese $\infty^{2h}$ selbst als eine $(\mu - r + h)$mal zählende besitze. Für $\mu + h \leqq r$ hat $\Theta$ keiner Bedingung zu genügen.

Eine eingehende Untersuchung des Verhaltens der vielfachen Integrale in allen hier möglichen Fällen in Bezug auf ihre Unstetigkeiten, sowie die noch weitergehende Untersuchung in Bezug auf ihre Werthe bei verschiedenen Integrationsgebieten, liegt unserm eigentlichen Gegenstande viel zu fern, als dass wir sie hier führen könnten. Beiläufig bemerke ich, dass schon die Doppelintegrale die Eigenschaft haben, bei stetig veränderten Integrationsgebieten \emph{zwischen denselben Grenzen} im Allgemeinen sich stetig ändernde Werthe zu geben${}^{*)}$.

\textbf{*)} Diese Untersuchungen führen unter andern interessanten Resultaten auch auf die Erweiterung des \textsc{Abel}'schen Theorems für vielfache Integrale, die schon \textsc{Jacobi} nach verschiedenen, an mehreren Orten zerstreuten Bemerkungen für speciellere Fälle gekannt hat. Dasselbe spricht sich für unsere Differentialausdrücke, in denen $\Theta$ eine ganze Function ist, folgendermassen aus:

„Die Summe
\[
\sum \frac{\Theta \cdot \Sigma \pm c_1\, x_2\, \frac{\partial x_3}{\partial \lambda_1} \frac{\partial x_4}{\partial \lambda_2} \dots \frac{\partial x_{r+2}}{\partial \lambda_r}}{\displaystyle\sum_i c_i\, \frac{\partial f}{\partial x_i}},
\]
ausgedehnt über die $n \cdot m_1 m_2 \dots m_r$ Schnittpunkte von $f = 0$ mit
\[
\varphi_1 + \lambda_1 \varphi'_1 = 0, \quad \varphi_2 + \lambda_2 \varphi'_2 = 0, \ \dots\ \varphi_r + \lambda_r \varphi'_r = 0,
\]
welche einem gegebenen Werthsysteme der $\lambda$ entsprechen, ist gleich Null.“

In speciellen Fällen, in denen $f$ mit einigen der $\varphi$ höhere Mannigfaltigkeiten gemein hat, reducirt sich natürlich die Anzahl der Glieder dieser Summe bedeutend.

Man beweist dieses Theorem entweder durch Betrachtungen, die den von \textsc{Riemann} über das einfache \textsc{Abel}'sche Theorem angestellten analog sind, oder rein algebraisch, indem man Zähler und Nenner des Ausdrucks $\Omega$ mit der Determinante
\[
\Sigma \pm k_1\, f_2 (\varphi_1 + \lambda_1 \varphi'_1)_3 \dots (\varphi_r + \lambda_r \varphi'_r)_{r+2},
\]
wo die $k_i$ Constanten und $(\varphi + \lambda \varphi')_i = \frac{\partial \varphi}{\partial x_i} + \lambda \frac{\partial \varphi'}{\partial x_i}$, multiplicirt und auf das Product den von \textsc{Jacobi} herrührenden, von \textsc{Clebsch} in Crelle's Journal Bd.~63, p.~228 bewiesenen Satz anwendet.

Für die nicht homogene Form $f(s,\, x_1,\, x_2,\, \dots\, x_r) = 0$ lautet der Satz:

„Die Summe:
\[
\sum \frac{Q}{\dfrac{\partial f}{\partial s}} \cdot \varDelta,
\]
wo $Q$ eine ganze Function $(n - r - 2)^{\text{ter}}$ Ordnung der $s$ und $x$, $\varDelta$ die Functionaldeterminante der $x$ nach den $\lambda$, ist, wenn ausgedehnt über die $n \cdot m_1 m_2 \dots m_r$ Schnittpunkte von $f = 0$ mit den $\varphi_i + \lambda_i \varphi'_i$, welche einem gegebenen Werthsysteme der $\lambda$ entsprechen, gleich Null.“

Die Gesammtheit der so entstehenden, den verschiedenen Werthen von $\Theta$ entsprechenden Gleichungen bildet ein simultanes System partieller Differentialgleichungen $r^{\text{ter}}$ Ordnung mit $r$ unabhängigen Variabeln, das eine particuläre algebraische Integration zulässt.

\origpage{305}

\section*{§ 4.}

\subsection*{Die eindeutigen Transformationen.}

Wir benutzen die vorstehenden Betrachtungen dazu, die Differentialausdrücke
\[
\Omega = \frac{\Theta \cdot \Sigma \pm c_1\, x_2\, \frac{\partial x_3}{\partial \lambda_1} \frac{\partial x_4}{\partial \lambda_2} \dots \frac{\partial x_{r+2}}{\partial \lambda_r}}{\displaystyle\sum_i c_i\, \frac{\partial f}{\partial x_i}}
\]
zu normiren. Wir legen der Function $\Theta$ die Eigenschaften bei, erstens,

\origpage{306}
eine \emph{ganze} Function $(n - r - 2)^{\text{ter}}$ Ordnung zu sein, und zweitens, jede $\mu$fach zählende $\infty^{2.h}$ auf $f = 0$ selbst als eine $(\mu - r + h)$fach zählende $\infty^{2.h}$ zu besitzen. Dabei sollen die in dem vorhergehenden §. für die Singularitäten der vielfachen $\infty^{2.h}$ angenommenen Voraussetzungen gelten. Wenn $\Theta$ diesen Forderungen gemäss bestimmt ist, so nennen wir der Kürze wegen das zugehörige $\Omega$ einen \emph{für $f = 0$ normirten Ausdruck $\Omega$}.

Die Anzahl $p$ der für $f = 0$ normirten Ausdrücke $\Omega$ ist eine endliche und gleich der Anzahl von Coefficienten, die in $\Theta$ noch unbestimmt bleiben, wenn $\Theta$ den obigen Bedingungen genügt, nämlich gleich der Anzahl solcher speciellen Functionen $\Theta'$, aus denen sich das allgemeinste, den Bedingungen genügende $\Theta$ linear zusammensetzt. Giebt es keine Mannigfaltigkeiten $\Theta$ mit den geforderten Eigenschaften, so ist die Anzahl $p = 0$ zu setzen.

Um unsern Zweck zu erreichen, eine charakteristische Zahl für die Irrationalität, die mit der Gleichung $f = 0$ verbunden ist, aufzustellen, werden wir die für $f = 0$ normirten Ausdrücke $\Omega$ in andere Ausdrücke eindeutig transformiren, die sich auf eine neue Gleichung $F(y) = 0$ beziehen, welche $f(x) = 0$ eindeutig entspricht. Es seien die rationalen Transformationsformeln:
\[
x_i = \varphi_i(y), \quad (i = 1,\, 2,\, \dots\, r+2),
\]
in denen die $\varphi$ ganze Functionen der $s^{\text{ten}}$ Ordnung. Durch diese Formeln geht $f(x) = 0$ über in:
\[
M \cdot F(y) = 0,
\]
wobei sich ein Factor $M$ nach §~1 absondern muss. Durch die gleiche Transformation wird sodann der für $f = 0$ normirte Ausdruck $\Omega$ übergehen in einen Ausdruck $\Omega'$:
\[
\Omega = \Omega' = \frac{\Theta \cdot D \cdot \Sigma \pm k_1\, y_2\, \frac{\partial y_3}{\partial \lambda_1} \frac{\partial y_4}{\partial \lambda_2} \dots \frac{\partial y_{r+2}}{\partial \lambda_r}}{s \cdot M \cdot \displaystyle\sum_i k_i\, \frac{\partial F}{\partial y_i}},
\]
wo
\[
\begin{gathered}
D = \Sigma \pm \frac{\partial x_1}{\partial y_1} \cdot \frac{\partial x_2}{\partial y_2} \dots \frac{\partial x_{r+2}}{\partial y_{r+2}}, \\
k_h = \sum_i c_i\, \frac{\partial y_h}{\partial x_i} .
\end{gathered}
\]

Auf den Werth von $\Omega'$ bleiben die Grössen $k_h$ ganz ohne Einfluss, und diese lassen sich daher als willkürliche Constanten, als die Coordinaten eines beliebigen Punktes, betrachten.

Der neue Ausdruck $\Omega'$ hat ganz die Form von $\Omega$. Wir werden nun nachweisen, dass, wenn das Entsprechen von $f(x) = 0$ und $F(y) = 0$ ein eindeutiges ist, der Zähler $\Theta D$ von $\Omega'$, der eine ganze Function

\origpage{307}
der $y$ ist, durch die gemeinschaftlichen Werthsysteme von $M = 0$ mit $F = 0$ ebenso oft als $M$ selbst hindurchgeht, woraus dann zu schliessen ist, dass $\frac{\Theta D}{M}$ mit Hülfe von $F = 0$ zu einer \emph{ganzen} Function gemacht werden kann; wir werden ferner zeigen, dass diese ganze Function $\frac{\Theta D}{M}$ in den vielfach zu zählenden Mannigfaltigkeiten von $F = 0$ diejenigen Eigenschaften besitzt, welche $\Omega'$ zu einem für $F = 0$ normirten Ausdrucke $\Omega$ machen. Bei eindeutigem Entsprechen kann man daher aus jedem für $f = 0$ normirten Ausdrucke $\Omega$ einen solchen in Bezug auf $F = 0$, und umgekehrt, ableiten, und die Anzahl $p$ dieser Ausdrücke ist folglich dieselbe in Bezug auf beide Gleichungen. Da sich diese Zahl $p$ aus dem Grade und den Singularitäten jeder Gleichung zusammensetzen lässt, so kann man sie als eine für die Irrationalität dieser Gleichung charakteristische Zahl aufstellen; ihre Gleichheit bei zwei $2r$fach unendlichen Mannigfaltigkeiten ist eine \emph{nothwendige} Bedingung für die Möglichkeit des eindeutigen Entsprechens derselben.

Die erwähnten Beweise werde ich hier für drei unabhängige Variable führen, aber in einer Form, die sich direct erweitern lässt und die Beweise allgemein gültig macht. Ich schicke einige Hülfssätze über das Verhalten der Substitutionsdeterminante $D$, die ich im Folgenden anzuwenden habe, voraus.

\section*{§ 5.}

\subsection*{Untersuchung der Determinante der Substitution.}

Durch die Transformationsformeln
\[
x_i = \varphi_i(y_1,\, y_2,\, y_3,\, y_4,\, y_5), \quad (i = 1,\, 2,\, \dots\, 5),
\]
wo die $\varphi$ von der $s^{\text{ten}}$ Ordnung, werde $f(x_1,\, x_2,\, \dots\, x_5) = 0$ in $F(y_1,\, y_2,\, \dots\, y_5) = 0$ übergeführt. Wir werden das Verhalten der Substitutionsdeterminante der Reihe nach in den Fällen untersuchen, in denen einer niedrigeren Mannigfaltigkeit auf $f = 0$ eine höhere auf $F = 0$, und umgekehrt, entspricht. Dabei können wir uns wieder auf die einfachsten Arten der Transformation beschränken, indem sich die übrigen Fälle auf diese mittelst des Princips der successiven Transformation zurückführen lassen.

1.~Wenn einem Punkte auf $f$ eine Fläche auf $F$ entspricht, so kann man die Transformationsformeln in die Form setzen:
\[
\begin{aligned}
\varrho x_1 &= A \cdot \psi_1 \\
\varrho x_2 &= A \cdot \psi_2 \\
\varrho x_3 &= A \cdot \psi_3 \\
\varrho x_4 &= A \cdot \psi_4 \\
\varrho x_5 &= \varphi_5 .
\end{aligned}
\]

\origpage{308}
Die Determinante
\[
D = \Sigma \pm \left(A \frac{\partial \psi_1}{\partial y_1} + \psi_1 \frac{\partial A}{\partial y_1}\right) \dots \left(A \frac{\partial \psi_4}{\partial y_4} + \psi_4 \frac{\partial A}{\partial y_4}\right) \cdot \frac{\partial \varphi_5}{\partial y_5}
\]
verschwindet im \emph{dritten} Grade längs des Schnittes von $A = 0$ mit $F = 0$.

2.~Wenn einem Punkte auf $f$ eine Curve auf $F$ entspricht, so werden die Transformationsformeln:
\[
\begin{aligned}
\varrho x_1 &= A \psi_1 + B \psi'_1 \\
\varrho x_2 &= A \psi_2 + B \psi'_2 \\
\varrho x_3 &= A \psi_3 + B \psi'_3 \\
\varrho x_4 &= A \psi_4 + B \psi'_4 \\
\varrho x_5 &= \varphi_5 .
\end{aligned}
\]
Die Transformationsdeterminante $D$ verschwindet hier im \emph{zweiten} Grade in dem Schnitt von $A = 0$, $B = 0$, $F = 0$.

3.~Wenn einer Curve auf $f$ eine Fläche auf $F$ entspricht, so kann man in der Nähe jedes Punktes dieser Fläche den Transformationsformeln die Gestalt 2. geben, und $D$ verschwindet daher ebenfalls \emph{zweimal} längs der auf $F$ liegenden Fläche.

In den umgekehrten Fällen verschwinden die 5 Functionen $\varphi_i$, jedoch, wie wir annehmen, nicht sämmtlich in höherm, als dem ersten Grade.

Es seien die Elemente von $D$ bezeichnet mit $\varphi_{ik} = \frac{\partial \varphi_i}{\partial y_k}$ und die Unterdeterminanten von $D$ seien $D_{ik} = \frac{\partial D}{\partial \varphi_{ik}}$. Wir gehen aus von der identischen Gleichung:
\[
D \cdot y_\lambda = s \cdot \sum_h D_{h\lambda} \cdot \varphi_h, \tag{$\alpha$}
\]
durch Differentiation nach $y_\mu$, wo $\mu$ von $\lambda$ verschieden, ergiebt sich:
\[
\frac{\partial D}{\partial y_\mu} \cdot y_\lambda = s \sum_h \left\{\frac{\partial D_{h\lambda}}{\partial y_\mu} \cdot \varphi_h + D_{h\lambda} \cdot \varphi_{h\mu}\right\};
\]
da aber identisch:
\[
\sum_h D_{h\lambda} \varphi_{h\mu} = 0,
\]
so folgt:
\[
\frac{\partial D}{\partial y_\mu} \cdot y_\lambda = s \sum_h \frac{\partial D_{h\lambda}}{\partial y_\mu} \cdot \varphi_h . \tag{$\beta$}
\]

4.~Die Functionen $\varphi$ sollen einen \emph{Punkt} gemein haben. In diesem Punkte verschwinden nach $(\alpha)$ und $(\beta)$ die Determinante $D$ und ihre ersten Differentialquotienten, d.~h. in einem gemeinsamen \emph{Schnittpunkte} der $\varphi$ verschwindet $D$ zweifach.

\origpage{309}
Eine weitere Differentiation nach $y_\nu$, wo $\nu$ von $\mu$ und $\lambda$ verschieden, ergiebt:
\[
\frac{\partial^2 D}{\partial y_\mu \partial y_\nu} \cdot y_\lambda = s \sum_h \frac{\partial D_{h\lambda}}{\partial y_\mu} \cdot \varphi_{h\nu} + s \sum_h \frac{\partial^2 D_{h\lambda}}{\partial y_\mu \partial y_\nu} \cdot \varphi_h . \tag{$\gamma$}
\]
Aber aus der identischen Gleichung $\displaystyle\sum_h D_{h\lambda} \varphi_{h\nu} = 0$ folgt die wiederum identische Gleichung:
\[
\sum_h \left\{\frac{\partial D_{h\lambda}}{\partial y_\mu} \cdot \varphi_{h\nu} + D_{h\lambda} \cdot \frac{\partial \varphi_{h\nu}}{\partial y_\mu}\right\} = 0, \tag{$\delta$}
\]
daher:
\[
\frac{\partial^2 D}{\partial y_\mu \partial y_\nu} \cdot y_\lambda = -\, s \sum_h D_{h\lambda} \cdot \frac{\partial \varphi_{h\nu}}{\partial y_\mu} + s \sum_h \frac{\partial^2 D_{h\lambda}}{\partial y_\mu \partial y_\nu} \cdot \varphi_h . \tag{$\gamma'$}
\]

5.~Die Functionen $\varphi$ sollen eine \emph{Curve} gemein haben. In diesem Falle kann man immer in jedem Punkte dieser Curve einen Parameter $q$ so bestimmen, dass nicht nur:
\[
\varphi_i = s \cdot \sum_k \varphi_{ik} \cdot y_k = 0, \quad (i = 1,\, 2,\, \dots\, 5),
\]
sondern auch in diesem Punkte:
\[
\sum_k \varphi_{ik} \cdot \frac{\partial y_k}{\partial q} = 0, \quad (i = 1,\, 2,\, \dots\, 5).
\]
Aus dem System dieser Gleichungen folgt sogleich, dass in diesem Punkte die sämmtlichen ersten Unterdeterminanten $D_{ik}$ von $D$ verschwinden. Daher verschwinden auch nach $(\gamma')$ die zweiten Differentialquotienten von $D$, d.~h. längs einer den $\varphi$ gemeinsamen Curve verschwindet die Determinante $D$ \emph{dreifach}.

Aus $(\gamma)$ folgt durch weitere Differentiation nach $y_\pi$, wo $\pi$ von $\lambda$, $\mu$, $\nu$ verschieden:
\[
\begin{aligned}
\frac{\partial^3 D}{\partial y_\mu \partial y_\nu \partial y_\pi} \cdot y_\lambda = s \sum_h &\left\{\frac{\partial^2 D_{h\lambda}}{\partial y_\mu \partial y_\pi} \cdot \varphi_{h\nu} + \frac{\partial D_{h\lambda}}{\partial y_\mu} \cdot \frac{\partial \varphi_{h\nu}}{\partial y_\pi} + \frac{\partial^3 D_{h\lambda}}{\partial y_\mu \partial y_\nu \partial y_\pi} \cdot \varphi_h \right.\\
&\quad \left.{}+ \frac{\partial^2 D_{h\lambda}}{\partial y_\mu \partial y_\nu} \cdot \varphi_{h\pi}\right\} .
\end{aligned}
\]
Die identische Gleichung $(\delta)$ giebt aber:
\[
\sum_h \left\{\frac{\partial^2 D_{h\lambda}}{\partial y_\mu \partial y_\pi} \cdot \varphi_{h\nu} + \frac{\partial D_{h\lambda}}{\partial y_\mu} \cdot \frac{\partial \varphi_{h\nu}}{\partial y_\pi} + \frac{\partial D_{h\lambda}}{\partial y_\pi} \cdot \frac{\partial \varphi_{h\nu}}{\partial y_\mu} + D_{h\lambda} \frac{\partial^2 \varphi_{h\nu}}{\partial y_\mu \partial y_\pi}\right\} = 0.
\]
Benutzt man die Beziehung, die aus dieser durch Vertauschung von $\nu$ und $\pi$ hervorgeht, so wird die vorhergehende Gleichung:
\[
\begin{aligned}
\frac{\partial^3 D}{\partial y_\mu \partial y_\nu \partial y_\pi} \cdot y_\lambda = s \sum_h &\left\{\frac{\partial^3 D_{h\lambda}}{\partial y_\mu \partial y_\nu \partial y_\pi} \cdot \varphi_h - \frac{\partial D_{h\lambda}}{\partial y_\pi} \cdot \frac{\partial \varphi_{h\nu}}{\partial y_\mu}\right. \\
&\quad{}- \left.\frac{\partial D_{h\lambda}}{\partial y_\mu} \cdot \frac{\partial \varphi_{h\pi}}{\partial y_\nu} - \frac{\partial D_{h\lambda}}{\partial y_\nu} \cdot \frac{\partial \varphi_{h\pi}}{\partial y_\mu} - D_{h\lambda} \cdot \frac{\partial^2 \varphi_{h\nu}}{\partial y_\mu \partial y_\pi} - D_{h\lambda} \cdot \frac{\partial^2 \varphi_{h\pi}}{\partial y_\mu \partial y_\nu}\right\} .
\end{aligned}
\tag{$\varepsilon$}
\]

\origpage{310}
6.~Die Functionen $\varphi$ sollen eine \emph{Fläche} gemein haben. Nach 5. ist $\varphi_h = 0$, $D_{h\lambda} = 0$. Ferner kann man hier in jedem Punkte der Fläche immer zwei Parameter $q$, $r$ so bestimmen, dass zugleich:
\[
\begin{gathered}
\varphi_i = s \sum_k \varphi_{ik} \, y_k = 0, \\
\sum_k \varphi_{ik} \, \frac{\partial y_k}{\partial q} = 0, \\
\sum_k \varphi_{ik} \, \frac{\partial y_k}{\partial r} = 0, \\
(i = 1,\, 2,\, \dots\, 5),
\end{gathered}
\]
woraus folgt, dass die sämmtlichen ersten und zweiten Unterdeterminanten $D_{ik}$ und $\frac{\partial D_{ik}}{\partial \varphi_{\lambda\mu}}$ von $D$ verschwinden. Da
\[
\frac{\partial D_{ik}}{\partial y_\pi} = \sum \frac{\partial D_{ik}}{\partial \varphi_{\lambda\mu}} \cdot \frac{\partial \varphi_{\lambda\mu}}{\partial y_\pi} ,
\]
so folgt auch, dass sämmtliche $\frac{\partial D_{ik}}{\partial y_\pi}$ verschwinden. Nach $(\varepsilon)$ verschwinden also auch die dritten Differentialquotienten von $D$; d.~h. längs einer den $\varphi$ gemeinsamen Fläche verschwindet die Determinante $D$ vierfach.

Hat man vier Functionen $\varphi$ von gleichem Grade, welche \emph{Flächen} bedeuten, so bezeichnet man die Functionaldeterminante dieser vier Functionen als die \textsc{Jacobi}'sche Fläche der $\varphi$, und die Sätze 4. und 5. bedeuten hier:

„In einem gemeinsamen Schnittpunkte der Flächen $\varphi_i = 0$ hat dio\ednote{Original reads: dio; presumably die intended.} \textsc{Jacobi}'sche Fläche derselben einen Doppelpunkt. Haben diese vier Flächen gleichen Grads eine Curve gemein, so ist diese Curve zugleich eine dreifache Curve der \textsc{Jacobi}'schen Fläche.“

\section*{§ 6.}

\subsection*{Untersuchung der transformirten Ausdrücke $\Omega$.}

Eine Untersuchung der Ausdrücke $\Omega'$ ist nur für diejenigen Punkte anzustellen, in denen $M$ oder in denen die $F'(y_h)$ sämmtlich verschwinden. Für diese Punkte hat man die Gleichungen:
\[
M \cdot F'(y_h) = 0 = \sum_i \frac{\partial f}{\partial x_i} \cdot \frac{\partial x_i}{\partial y_h}, \quad (h = 1,\, 2,\, \dots\, 5).
\]
Mittelst dieser Gleichungen kann man schon in den einfacheren Fällen, in denen niedrigeren Mannigfaltigkeiten auf $f$ höhere auf $F$ entsprechen, zeigen, dass der Zähler $\Theta \cdot D$ in $\Omega'$ durch die $M$ und

\origpage{311}
$F$ gemeinschaftlichen Werthsysteme wenigstens ebenso oft als $M$ hindurchgeht.

1.~Einem \emph{einfachen} Punkte auf $f$ entspreche eine Fläche auf $F$. Nach den Transformationsformeln §~5.~1. wird sich $M$ hier wie $A$, die Determinante $D$ wie $A^3$ verhalten. $\frac{D}{M}$ muss also längs der Fläche \emph{zweimal} verschwinden, und ebenso $\Omega'$. Diesen Umstand werden wir noch weiterhin zu benutzen haben.

2.~Einem \emph{Doppelpunkt} auf $f$ entspreche eine Fläche auf $F$. Nach §~5.~1. verhält sich $M$ wie $A^2$, $D$ wie $A^3$, d.~h. $\frac{D}{M}$ verschwindet \emph{einmal} längs der Fläche.

3.~Einem höhern Knotenpunkte auf $f$ entspreche eine Fläche auf $F$. Man hat hier, wie schon bei 2.:
\[
\frac{\partial f}{\partial x_1} = 0, \quad \frac{\partial f}{\partial x_2} = 0, \quad \dots \quad \frac{\partial f}{\partial x_5} = 0.
\]
$\Theta$ verschwindet dabei, wenn der Knotenpunkt ein $\mu$facher ist, $(\mu - 3)\text{mal}$, $M$ verhält sich wie $A^\mu$, $D$ wie $A^3$. $\Theta D$ verschwindet also ebenso oft wie $M$.

4.~Einer einfachen Curve auf $f$ entspreche eine Fläche auf $F$. Nach §~5.~3. verschwindet $D$ zweimal, während $M$ einmal verschwinden kann, also $D$ jedenfalls einmal mehr als $M$, längs der Fläche, welchen Umstand wir ebenfalls noch zu benutzen haben werden.

5.~Einer $\mu$fachen Curve auf $f$ entspreche eine Fläche auf $F$. Hier verschwindet $\Theta$ $(\mu - 2)\text{mal}$, $D$ zweimal, $M$ kann $\mu\text{mal}$ verschwinden, also $\Theta D$ ebenso oft als $M$.

6.~Einer $\mu$fachen Fläche entspreche eine Fläche auf $F$. Da die $\frac{\partial f}{\partial x_i}$ $(\mu - 1)\text{mal}$ verschwinden, muss nach der obigen Gleichung auch $M$ $(\mu - 1)\text{mal}$ verschwinden; aber $\Theta$ verschwindet hier ebenso oft als $M$.

Die Untersuchungen der umgekehrten Fälle, in denen die $\varphi$ sämmtlich einfach, also die $\frac{\partial f}{\partial x_i}$ und ebenso die $M F'(y_h)$ $(n - 1)\text{fach}$ verschwinden, erfordert eine nähere Betrachtung der Abbildungen der auf $f = 0$ liegenden vielfachen Mannigfaltigkeiten. Ich werde hierbei, mit den entsprechenden Erweiterungen, den Weg verfolgen, der von \textsc{Clebsch} und \textsc{Gordan}, \textsc{Abel}'sche Functionen, §~16., für Curven und von \textsc{Clebsch}, diese Zeitschrift, Bd.~I, p.~270, für die Abbildung der Doppelcurven rationaler Flächen eingeschlagen worden ist.

Wenn $y$ und $y'$ zwei Werthsysteme sind, denen ein Werthsystem der $x$ entspricht, so hat man, um die Abbildung der auf $f = 0$ liegenden vielfachen Mannigfaltigkeiten zu finden, aus den Bedingungen:

\origpage{312}
\[
\begin{cases}
\varphi_i(y) = \varphi_i(y') \\
F(y') = 0
\end{cases}
\tag{1}
\]
die $y'$ zu eliminiren. Um nicht auf identische Gleichungen zu kommen, lässt man den Punkt $x' = \varphi(y')$ sich zunächst dem Punkt $x$ nur annähern, indem man
\[
x'_i = \varphi_i(y') = \varphi_i(y) + \varepsilon z_i \tag{2}
\]
setzt, und eliminirt aus (2) und $F(y')$ die $y'$. Wenn man noch in der Resultante dieser Gleichungen $\varepsilon$ gegen Null convergiren lässt, schreibt sich dieselbe in der Form:
\[
0 = R = z_1 f'(\varphi_1) + z_2 f'(\varphi_2) + \dots + z_5 f'(\varphi_5) . \tag{3}
\]
Das Schnittsystem von $R = 0$, $F(y) = 0$ hängt noch von den ganz willkürlichen Grössen $z$ ab, die durch die Art der Elimination (2) eingehen und in einem unwesentlichen Factor herauszuschaffen sind. Die Gleichungen (2) werden nun noch erfüllt, wenn sich die $y'$ von den $y$ nur um gewisse Grössen von der Ordnung $\varepsilon$ unterscheiden, wenn man also setzt:
\[
y'_i = y_i + \varepsilon \eta_i ,
\]
und wenn zugleich die $y$ einer weitern von den $z$ abhängigen Gleichung genügen, die sich durch Elimination aus den Gleichungen ergiebt, in die jetzt die (2) übergehen:
\[
\begin{aligned}
\frac{\partial \varphi_i(y)}{\partial y_1} \cdot \eta_1 + \dots + \frac{\partial \varphi_i(y)}{\partial y_5} \eta_5 &= z_i , \\
\frac{\partial F}{\partial y_1} \cdot \eta_1 + \dots + \frac{\partial F}{\partial y_5} \eta_5 &= 0 ,
\end{aligned}
\]
nämlich:
\[
S = 0 =
\begin{vmatrix}
\varphi'_1(y_1), & \varphi'_1(y_2), & \dots & \varphi'_1(y_5), & z_1 \\
\cdot & \cdot & \dots & \cdot & \cdot \\
\cdot & \cdot & \dots & \cdot & \cdot \\
\varphi'_5(y_1), & \varphi'_5(y_2), & \dots & \varphi'_5(y_5), & z_5 \\
F'(y_1), & F'(y_2), & \dots & F'(y_5), & 0
\end{vmatrix} . \tag{4}
\]
So lange die Coefficienten der $z$ keinen gemeinschaftlichen Factor haben, werden die auf $f = 0$ liegenden vielfachen Gebilde durch die Gleichungen:
\[
\frac{R}{S} = 0, \quad F(y) = 0
\]
abgebildet. Da hier jedoch im Allgemeinen auch auf $f = 0$ Fundamental-Punkte und -Curven liegen, so wird $S$ in Factoren zerfallen, von denen einige nicht in $R$ enthalten sein müssen. Wir haben hier diese Fälle genauer zu discutiren.

$R = 0$ geht, nach dem Ausdrucke (3), $(\mu - 1)\text{mal}$ durch die Abbildung jeder $\mu$fach zählenden, auf $f = 0$ liegenden Mannigfaltigkeit hindurch. Um das entsprechende Verhalten von $S$ zu finden, könnte

\origpage{313}\ednote{The running header of this page in the original print is misnumbered 213.}
man die Coefficienten der $z$ in (4) auf dem nämlichen Wege untersuchen, wie die Determinante $D$ in §~5. behandelt wurde. Wir können indess auch aus dem Verhalten von $D$ auf das von $S$ schliessen. Man hat nämlich, da die $z_i$ beliebige Grössen sind:
\[
\frac{R}{S} = \frac{f'(\varphi_1)}{\Sigma \pm \varphi'_2(y_1) \dots \varphi'_5(y_4) F'(y_5)} ,
\]
und gleich den analogen Quotienten. Aus den Gleichungen:
\[
M F'(y_h) = \sum_i f'(\varphi_i) \cdot \varphi'_i(y_h)
\]
folgt aber, dass diese Quotienten auch gleich $\frac{M}{D}$, also:
\[
\frac{R}{S} = \frac{M}{D} .
\]
$D$ verschwindet längs einer Fläche auf $F$, welche einem einfachen und einem Doppelpunkte auf $f$ entspricht, resp. zweimal und einmal mehr als $M$, und längs einer Fläche, welche einer einfachen Curve auf $f$ entspricht, einmal mehr als $M$. Die entsprechenden Factoren, welche auch in $S$ enthalten sein müssen, sind nicht in $R$ enthalten. Mit Hülfe der am Anfang dieses § für $\frac{M}{D}$ erhaltenen Resultate sieht man überhaupt, dass, wenn einem Punkte oder einer Curve auf $f$ eine Fläche auf $F$ entspricht, $S$ den entsprechenden Factor resp. zweimal oder einmal enthält. Wir sondern nun von $S$ diejenigen Factoren doppelt ab, welche einfachen Punkten auf $f$, und diejenigen einfach ab, welche Doppelpunkten und einfachen Fundamentalcurven auf $f$ entsprechen, und bezeichnen den Rest mit $S'$. Dieses $S'$ ist in $R$ mit Hülfe von $F = 0$ theilbar. $\frac{R}{S'}$ geht noch $(\mu - 1)\text{fach}$ durch die Abbildung der $\mu$fachen Fläche, $(\mu - 2)\text{fach}$ durch die einer $\mu$fachen Curve, und $(\mu - 3)\text{fach}$ durch die eines $\mu$fachen Punktes hindurch. Die Function $\Theta$ geht nun genau ebenso oft durch diese gemeinschaftlichen Werthsysteme von $\frac{R}{S'}$ und $F$ hindurch, folglich, da $\frac{R}{S} = \frac{M}{D}$, auch ebenso oft wie $\frac{M}{D'}$, wo $D'$ aus $D$ gerade so entsteht, wie $S'$ aus $S$.

Hiermit ist bewiesen, dass $\Theta$ durch alle Punkte, in denen die ganze Function $\frac{M}{D'}$ verschwindet, ebenso oft als $\frac{M}{D'}$, also wenigstens ebenso oft als $\frac{M}{D}$ hindurchgeht.

7.~Wir können jetzt mit wenigen Worten alle die Fälle erledigen, in denen höheren Mannigfaltigkeiten auf $f$ niedrigere auf $F$ entsprechen. Hier verschwinden die sämmtlichen $x$ einfach, daher $\Theta$, als Function der $(n - 5)^{\text{ten}}$ Ordnung, im Allgemeinen $(n - 5)\text{mal}$. Aber nach dem eben bewiesenen Satze tritt hier, vermöge der auf $f = 0$

\origpage{314}
liegenden singulären Mannigfaltigkeiten, der besondere Umstand ein, dass $\Theta$ in den einfachen Fundamentalgebilden von $F = 0$ in ebenso hoher Ordnung als $\frac{M}{D}$, also in höherer Ordnung, als der $(n - 5)^{\text{ten}}$ verschwindet. So in einem Punkte, der einer Fläche auf $f$ entspricht, $(n - 3)\text{mal}$, in einer einfachen Fundamentalcurve $(n - 4)\text{mal}$, und ebenso oft in einem Doppelpunkte, der einer Fläche auf $f = 0$ entspricht (da hier $M$ $(n - 2)\text{mal}$, $D$ zweimal verschwindet). Man sieht, dass, wenn sich zwei Flächen eindeutig entsprechen, die Ausdrücke $\Omega$ längs derjenigen Mannigfaltigkeiten auf der einen Fläche, welche einfachen Fundamentalgebilden der andern Fläche entsprechen, immer solche Singularitäten zeigen müssen, wie sie unter 1. und 4. dieses § erwähnt sind.

8.~Wenn auf $F$ ein $\mu$facher Knotenpunkt entsteht, werden die $M F'(y_h)$ Null in der $(n - 1)^{\text{ten}}$ Ordnung, also $M$ in der $(n - \mu)^{\text{ten}}$ Ordnung. $D$ verschwindet nach §~5.~4. zweifach; $\Theta$ verschwindet $(n - 5)\text{fach}$, also $\frac{\Theta D}{M}$ $(\mu - 3)\text{fach}$. Wenn auf $F$ eine $\mu$fache Curve entsteht, findet man ebenso mit Hülfe von §~5.~5., dass $\frac{\Theta D}{M}$ $(\mu - 2)\text{fach}$ verschwindet, und wenn eine $\mu$fache Fläche entsteht, mit Hülfe von §~5.~6., dass $\frac{\Theta D}{M}$ $(\mu - 1)\text{fach}$ verschwindet.

Strenge genommen, haben wir sub 1. bis 7. nur nachgewiesen, dass $\Theta D$ durch das ganze Schnittsystem von $M = 0$, $F = 0$ hindurchgeht, und es müsste noch der in der Curven- und Flächentheorie fortwährend gebrauchte Satz dargethan werden, dass $\Theta D$ sich deshalb nothwendiger Weise in die Form setzen lässt:
\[
\Theta D \equiv A \cdot M + B \cdot F .
\]
Es seien $\Theta D$, $M$, $F$ Functionen der $s$ und $x$, $F$ vom Grade $m$. Man kann nun $M$ immer mit einem solchen Factor $\lambda$ multipliciren, dass aus $M \cdot \lambda$ mit Hülfe von $F = 0$ alles $s$ verschwindet, und hat dann:
\[
M \lambda = \Phi + \mu \cdot F ,
\]
wo $\Phi$ nur noch eine Function der $x$. Dividirt man nun $\lambda \cdot \Theta D$ durch $F$, so lässt sich der Rest auf die $(m - 1)^{\text{te}}$ Potenz in Bezug auf $s$ erniedrigen, also:
\[
\Theta D \cdot \lambda = \nu \cdot F + \Phi_1 s^{m-1} + \Phi_2 \cdot s^{m-2} + \dots + \Phi_m ,
\]
wo die $\Phi$ Functionen der $x$. Für jedes Werthsystem der $x$, das $\Phi$ zu Null macht, lassen sich $m$ Wurzeln für $s$ angeben, für welche auch $F$ verschwindet. Für alle diese Werthsysteme verschwindet aber auch $M$ oder $\lambda$, also $\Theta D \cdot \lambda$, daher auch der Ausdruck:
\[
\Phi_1 s^{m-1} + \Phi_2 s^{m-2} + \dots + \Phi_m .
\]

\origpage{315}
Da nun dieser Ausdruck nur für $(m - 1)$ Wurzeln zu $0$ werden könnte, müssen alle $\Phi_i$ entweder identisch verschwinden, oder $\Phi$ zum Factor haben, und man hat die Darstellung:
\[
\Theta D \cdot \lambda = \nu F + \varrho \Phi = (\nu - \varrho \mu) F + \varrho \lambda \cdot M ,
\]
oder:
\[
\Theta D \equiv \varrho M + \frac{\nu - \varrho \mu}{\lambda} \cdot F .
\]
Da $\Theta D$ eine ganze Function ist, so muss, wenn $F$ und $\lambda$ keinen gemeinschaftlichen Factor haben, $\lambda$ in $\nu - \varrho \mu$ theilbar sein. Ein gemeinschaftlicher Factor von $\lambda$ und $F$ könnte nach der Gleichung $M \lambda = \Phi + \mu F$ nur von den $x$ abhängig sein, ein Fall, der durch eine vorhergehende lineare Transformation immer zu vermeiden ist. Die verlangte Darstellung ist somit für jeden Fall bewiesen.

\section*{§ 7.}

\subsection*{Das Geschlecht der algebraischen Gebilde.}

Die bisher geführten Untersuchungen liefern den Beweis, dass den für $f = 0$ normirten Ausdrücken $\Omega$ bei eindeutiger Transformation wieder für $F = 0$ normirte Ausdrücke entsprechen, und dass also deren Anzahl $p$ die gleiche ist. Um das Gesammtgebiet der hier behandelten Fälle zu übersehen, braucht man nur wieder successive Transformationen auf die speciell betrachteten Fälle anzuwenden, um zu dem Resultate zu kommen, dass alle Gleichungen, die $\infty^{2.r}$ mit vielfachen $\infty^{2.h}$ ausdrücken, welche nicht selbst wieder solche specielle Singularitäten besitzen, wie sie in §~3. bezeichnet sind, in die Betrachtung eingeschlossen sind. Wir bezeichnen die Zahl $p$ als das \emph{Geschlecht} der $\infty^{2.r}$, $f = 0$, und haben, um das ganze Gebiet der algebraischen Gleichungen mit den eben erwähnten Beschränkungen nach dieser Zahl $p$ in Classen einzutheilen, in diesem Umfange das folgende Theorem bewiesen:

Das Geschlecht $p$ der $\infty^{2.r}$, $f = 0$, ist gleich der Anzahl der in einer $\infty^{2.r}$ vom Grade $n - r - 2$, $\Theta = 0$, noch unbestimmt bleibenden Constanten, wenn $\Theta = 0$ gezwungen wird, durch jede $\mu$fache auf $f = 0$ liegende $\infty^{2.h}$ $(\mu - r + h)\text{fach}$ hindurchzugehen. Für $\mu + h \leqq r$ hat $\Theta$ keiner Bedingung zu genügen.
Die Zahl $p$ hat die Eigenschaft, für zwei $\infty^{2.r}$, $f = 0$ und $F = 0$, welche sich im Allgemeinen Punkt für Punkt eindeutig entsprechen, die gleiche zu sein, so dass $p$ die Classenzahl ist für die Classe, zu der $f = 0$ und $F = 0$ gehören. Die Gleichheit der Zahl $p$ ist das \emph{nothwendige} Kriterium für die Möglichkeit

\origpage{316}
einer eindeutigen rationalen Transformation zweier $\infty^{2.r}$ in einander.

Die Frage nach den weiteren und hinreichenden Bedingungen des eindeutigen Entsprechens, die ich schon in meiner oben citirten Notiz behandelt habe, muss ferneren Untersuchungen vorbehalten bleiben.

Göttingen, den 16.~August 1869.

\end{document}
